\documentclass[12pt]{article}
\usepackage{fullpage,amsmath,amssymb,amsthm,hyperref,amsfonts}
\newtheorem{theorem}{Theorem}
\newtheorem{lemma}[theorem]{Lemma}
\newtheorem{proposition}[theorem]{Proposition}
\newtheorem{corollary}[theorem]{Corollary}
\theoremstyle{definition}
\newtheorem{definition}[theorem]{Definition}
\newtheorem{remark}[theorem]{Remark}

\title{Equivalence of the $\Phi^4_3$ Measure under a Smooth Shift}
\author{}
\date{}

\begin{document}
\maketitle

\section*{Statement and Answer}

Let $\mathbb{T}^3 = (\mathbb{R}/\mathbb{Z})^3$ be the unit three-torus, let $\mu$ denote the
$\Phi^4_3$ Euclidean field measure on $\mathcal{D}'(\mathbb{T}^3)$, let
$\psi:\mathbb{T}^3\to\mathbb{R}$ be a smooth function (not identically zero), and let
$T_\psi:\mathcal{D}'(\mathbb{T}^3)\to\mathcal{D}'(\mathbb{T}^3)$ be the shift
$T_\psi(u)=u+\psi$. Write $T_\psi^*\mu$ for the pushforward of $\mu$ by $T_\psi$.

\medskip

\noindent\textbf{Answer.} \emph{Yes:} the measures $\mu$ and $T_\psi^*\mu$ are equivalent
(mutually absolutely continuous), with the explicit Radon--Nikodym derivative computed in
Theorem~\ref{thm:final}.

\medskip

\noindent\textbf{Strategy.} Write $\mu$ as a Radon--Nikodym density times the underlying
massive Gaussian free field $\mu_0$. Since $\psi$ is smooth, it lies in the Cameron--Martin
space of $\mu_0$, so the Cameron--Martin theorem gives $T_\psi^*\mu_0\sim\mu_0$ with an
explicit density. The renormalised $\Phi^4$ density transforms under the shift by a
polynomial expression in Wick powers of $\phi$ of total chaos degree at most three with
\emph{smooth} coefficients depending on $\psi$. The product of the Cameron--Martin density
and the exponentiated Wick-polynomial difference is shown to be a positive,
$\mu_0$-integrable random variable, hence a valid Radon--Nikodym derivative.

\section{Notation and Construction of $\mu$}

\subsection{The free field}

Fix $m^2>0$; the precise value is irrelevant, only that $-\Delta+m^2$ is invertible on
$\mathcal{D}'(\mathbb{T}^3)$. Let $C := (-\Delta+m^2)^{-1}$ acting on $L^2(\mathbb{T}^3)$.
Diagonalising in Fourier,
\begin{equation}\label{eq:FourierC}
\widehat{Cf}(k) \;=\; (4\pi^2|k|^2+m^2)^{-1}\,\widehat{f}(k),\qquad k\in\mathbb{Z}^3,
\end{equation}
so $C$ is a positive, self-adjoint, Hilbert--Schmidt-regularised pseudodifferential
operator of order $-2$.

\begin{definition}[Massive Gaussian Free Field]\label{def:GFF}
The \emph{massive Gaussian free field} on $\mathbb{T}^3$ is the centred Gaussian probability
measure $\mu_0$ on $\mathcal{D}'(\mathbb{T}^3)$ whose covariance is $C$:
\begin{equation}\label{eq:GFFcov}
\int_{\mathcal{D}'} \langle\phi,f\rangle\langle\phi,g\rangle\,d\mu_0(\phi) \;=\; \langle f, Cg\rangle_{L^2},\qquad f,g\in C^\infty(\mathbb{T}^3).
\end{equation}
\end{definition}

The Cameron--Martin space of $\mu_0$ is
\begin{equation}\label{eq:CMspace}
\mathcal{H} \;:=\; C^{1/2}\bigl(L^2(\mathbb{T}^3)\bigr)
\;=\; H^1(\mathbb{T}^3),
\end{equation}
the standard Sobolev space, with inner product
$\langle h,h'\rangle_{\mathcal{H}}=\langle h,(-\Delta+m^2)h'\rangle_{L^2}$. A direct Fourier
computation, using \eqref{eq:FourierC}, shows that $\mu_0$-a.e.\ realisation $\phi$ lies in
the Besov space $\mathcal{C}^{-1/2-\kappa}(\mathbb{T}^3)$ for every $\kappa>0$, but
$\phi\notin L^2(\mathbb{T}^3)$ almost surely. In particular, $\mu_0$ is genuinely supported
on a space of distributions, not functions.

\subsection{The renormalised potential}

For $N\in\mathbb{N}$ let $\Pi_N$ denote the Fourier projector onto frequencies
$\{|k|_\infty\le N\}$ and write $\phi_N := \Pi_N\phi$, which is a smooth function for
$\mu_0$-a.e.\ $\phi$. Define the variance
\begin{equation}\label{eq:CN}
c_N \;:=\; \mathbb{E}_{\mu_0}[\phi_N(x)^2] \;=\; \sum_{|k|_\infty\le N,\,k\ne0}\frac{1}{4\pi^2|k|^2+m^2} \;+\; \frac{1}{m^2},
\end{equation}
which is independent of $x\in\mathbb{T}^3$ by translation invariance and grows like
$c_N\sim CN$ as $N\to\infty$. Let
\begin{equation}\label{eq:Wick}
\colon\!\phi_N^2\!\colon\;:=\;\phi_N^2-c_N,\qquad
\colon\!\phi_N^3\!\colon\;:=\;\phi_N^3-3c_N\phi_N,\qquad
\colon\!\phi_N^4\!\colon\;:=\;\phi_N^4-6c_N\phi_N^2+3c_N^2,
\end{equation}
the Wick-renormalised powers with respect to the variance $c_N$
(Hermite-polynomial normalisation: with $H_n$ the probabilists' Hermite polynomials
satisfying the generating identity $\sum_n t^n H_n(x;c)/n!=\exp(tx-t^2c/2)$, one has
$\colon\!\phi_N^n\!\colon\,=H_n(\phi_N;c_N)$).

The $\Phi^4_3$ measure is the limit, as $N\to\infty$, of the regularised tilted measures
\begin{equation}\label{eq:muN}
d\mu_N \;:=\; \frac{e^{-V_N(\phi)}}{Z_N}\,d\mu_0(\phi),\qquad
V_N(\phi) \;:=\; \int_{\mathbb{T}^3}\Bigl(\tfrac{\lambda}{4}\colon\!\phi_N^4\!\colon
+ a_N\colon\!\phi_N^2\!\colon + b_N\Bigr)\,dx,
\end{equation}
with coupling $\lambda>0$ and renormalisation constants $a_N,b_N$ chosen so that the
limit exists.

\begin{theorem}[Existence of $\Phi^4_3$, Glimm--Jaffe]\label{thm:GJ}
There exist $a_N,b_N\in\mathbb{R}$ and $\lambda>0$ such that the measures $\mu_N$ in
\eqref{eq:muN} converge weakly on $\mathcal{D}'(\mathbb{T}^3)$ to a probability measure
$\mu$ as $N\to\infty$. The limit measure $\mu$ is supported on
$\mathcal{C}^{-1/2-\kappa}(\mathbb{T}^3)$ for every $\kappa>0$, and is absolutely continuous
with respect to $\mu_0$ with strictly positive density
\begin{equation}\label{eq:density}
\frac{d\mu}{d\mu_0}(\phi) \;=\; \frac{e^{-V(\phi)}}{Z},\qquad
V(\phi) := \lim_{N\to\infty} V_N(\phi)\ \text{in}\ L^p(\mu_0)\ \text{for every}\ p\in[1,\infty),
\end{equation}
with $0<Z<\infty$ and $e^{\pm V}\in L^p(\mu_0)$ for every $p\in[1,\infty)$.
\end{theorem}

\begin{proof}[Reference]
This is Theorem~5.6.1 in Glimm--Jaffe \cite{GlimmJaffe}, combined with the convergence
statement in their Section~23.1; for a more recent treatment via paracontrolled
distributions see Catellier--Chouk \cite{CatellierChouk} and Gubinelli--Hofmanov\'a
\cite{GH}, and via regularity structures Hairer \cite{HairerRS}. The two-sided
$L^p$-bound on $e^{\pm V}$ follows from the standard Wick-power $L^p$ estimates and the
inequality $V\ge -K_p$ ($\mu_0$-a.s.) controlled by $\colon\!\phi^4\!\colon$, also given in
\cite[Ch.~8]{GlimmJaffe}.
\end{proof}

We adopt \eqref{eq:density} as the definition of $\mu$. Two facts will be used repeatedly:
\begin{itemize}
\item[(F1)] $V:\mathcal{D}'(\mathbb{T}^3)\to\mathbb{R}\cup\{+\infty\}$ is a measurable
functional, finite $\mu_0$-a.e., with $V\in L^p(\mu_0)$ and $e^{\pm V}\in L^p(\mu_0)$ for
every $p\in[1,\infty)$ (Theorem~\ref{thm:GJ});
\item[(F2)] $V$ is a polynomial in the Wick-renormalised distributions
$\colon\!\phi^n\!\colon$, $n=2,3,4$, with constant coefficients (after the renormalising
counterterms have been absorbed into the constants $a_N,b_N$ in \eqref{eq:muN}); precisely,
\begin{equation}\label{eq:V}
V(\phi) \;=\; \frac{\lambda}{4}\int_{\mathbb{T}^3}\!\colon\!\phi^4\!\colon dx
\;+\; a\int_{\mathbb{T}^3}\!\colon\!\phi^2\!\colon dx \;+\; b
\end{equation}
in the sense of $L^p(\mu_0)$ limits; here $a,b\in\mathbb{R}$ are the renormalised
coefficients.
\end{itemize}

\section{Cameron--Martin for the Free Field}

\begin{theorem}[Cameron--Martin; Bogachev \cite{Bogachev}, Theorem 2.4.5]\label{thm:CM}
Let $\nu$ be a centred Gaussian Radon measure on a separable locally convex space $X$ with
Cameron--Martin space $\mathcal{H}\subseteq X$. For $h\in X$, the shifted measure $\nu_h$
defined by $\nu_h(A):=\nu(A-h)$ (i.e.\ the pushforward of $\nu$ under $x\mapsto x+h$) is
\begin{itemize}
\item equivalent to $\nu$ if $h\in\mathcal{H}$, with Radon--Nikodym derivative
\begin{equation}\label{eq:RNCM}
\frac{d\nu_h}{d\nu}(x) \;=\; \exp\!\left(\hat h(x)-\tfrac{1}{2}\|h\|_{\mathcal{H}}^2\right),
\end{equation}
where $\hat h\in L^2(\nu)$ is the unique centred Gaussian variable on $X$ such that
$\int x^*(x)\hat h(x)\,d\nu(x)=\langle x^*,h\rangle$ for every $x^*\in X^*$;
\item singular with respect to $\nu$ if $h\notin\mathcal{H}$.
\end{itemize}
\end{theorem}

We apply this to $\nu=\mu_0$ on $X=\mathcal{D}'(\mathbb{T}^3)$ and $h=\psi$.

\begin{lemma}\label{lem:CMfree}
For every $\psi\in C^\infty(\mathbb{T}^3)$ the measures $\mu_0$ and $T_\psi^*\mu_0$ are
equivalent and
\begin{equation}\label{eq:RPsi}
R_\psi(\phi) \;:=\; \frac{d(T_\psi^*\mu_0)}{d\mu_0}(\phi)
\;=\; \exp\!\Bigl(\langle\phi,(-\Delta+m^2)\psi\rangle - \tfrac12\langle\psi,(-\Delta+m^2)\psi\rangle\Bigr),
\end{equation}
where the pairing $\langle\phi,(-\Delta+m^2)\psi\rangle$ is interpreted as the
$L^2(\mu_0)$ centred Gaussian random variable with variance
$\langle\psi,(-\Delta+m^2)\psi\rangle$.
\end{lemma}

\begin{proof}
We verify the hypotheses of Theorem~\ref{thm:CM}. The space $X=\mathcal{D}'(\mathbb{T}^3)$
is separable (because the test-function space $\mathcal{D}(\mathbb{T}^3)=C^\infty(\mathbb{T}^3)$
is a separable Fr\'echet space) and locally convex; $\mu_0$ is a centred Gaussian measure
which is Radon on this Polish space \cite[II.5]{Bogachev}. Its Cameron--Martin space is
$\mathcal{H}=H^1(\mathbb{T}^3)$ by \eqref{eq:CMspace}, and $\psi\in
C^\infty(\mathbb{T}^3)\subset H^1(\mathbb{T}^3)$ since $\mathbb{T}^3$ is compact. Hence the
first alternative of Theorem~\ref{thm:CM} applies, giving
\begin{equation}\label{eq:RPsi-prelim}
\frac{d(T_\psi^*\mu_0)}{d\mu_0}(\phi) \;=\; \exp\!\Bigl(\hat\psi(\phi) - \tfrac12\|\psi\|_{\mathcal{H}}^2\Bigr).
\end{equation}
We identify $\hat\psi$ explicitly. By the abstract definition,
$\int\langle\phi,f\rangle\hat\psi(\phi)\,d\mu_0(\phi)=\langle f,\psi\rangle_{L^2}$ for every
test function $f\in C^\infty(\mathbb{T}^3)$. Since
$\langle f,\psi\rangle_{L^2}=\langle f,C(-\Delta+m^2)\psi\rangle_{L^2}$ (because
$C(-\Delta+m^2)=I$ on $L^2$) and the latter equals
$\mathbb{E}_{\mu_0}\!\bigl[\langle\phi,f\rangle\langle\phi,(-\Delta+m^2)\psi\rangle\bigr]$ by
\eqref{eq:GFFcov}, the centred Gaussian random variable $\hat\psi$ is given by
$\hat\psi(\phi)=\langle\phi,(-\Delta+m^2)\psi\rangle$ in $L^2(\mu_0)$. Moreover
$\|\psi\|_{\mathcal{H}}^2=\langle\psi,(-\Delta+m^2)\psi\rangle$ from \eqref{eq:CMspace}.
Substituting into \eqref{eq:RPsi-prelim} gives \eqref{eq:RPsi}.
\end{proof}

\begin{lemma}\label{lem:Lp-Rpsi}
For every $\psi\in C^\infty(\mathbb{T}^3)$ and every $\alpha\in\mathbb{R}$, $p\in[1,\infty)$,
$R_\psi^\alpha\in L^p(\mu_0)$.
\end{lemma}

\begin{proof}
By \eqref{eq:RPsi}, $R_\psi^\alpha(\phi)=\exp(\alpha X-\tfrac{\alpha}{2}\sigma^2)$ where
$X:=\langle\phi,(-\Delta+m^2)\psi\rangle$ is centred Gaussian with variance
$\sigma^2:=\langle\psi,(-\Delta+m^2)\psi\rangle<\infty$ (smoothness of $\psi$). Then
$R_\psi^{p\alpha}$ has expectation
$\mathbb{E}[\exp(p\alpha X-\tfrac{p\alpha}{2}\sigma^2)]=\exp(\tfrac{(p\alpha)^2-p\alpha}{2}\sigma^2)<\infty$
by the Gaussian Laplace transform.
\end{proof}

\section{Behaviour of Wick Powers under a Smooth Shift}

We now compute $V(\phi+\eta)$ in terms of Wick polynomials of $\phi$ for any smooth
deterministic shift $\eta\in C^\infty(\mathbb{T}^3)$. The key ingredient is the binomial
identity for Hermite polynomials under shift, applied componentwise.

\begin{lemma}[Wick shift formula]\label{lem:Wickshift}
Let $\phi$ be the GFF and $\eta\in C^\infty(\mathbb{T}^3)$. For each $n\in\mathbb{N}$, the
Wick power $\colon\!(\phi+\eta)^n\!\colon$ (renormalised by the variance $c_N$ of $\phi_N$,
not by that of $\phi_N+\eta$) satisfies
\begin{equation}\label{eq:Wickbin}
\colon\!(\phi+\eta)^n\!\colon \;=\; \sum_{k=0}^n\binom{n}{k}\eta^{n-k}\colon\!\phi^k\!\colon
\end{equation}
in the sense of $L^p(\mu_0)$ for every $p\in[1,\infty)$, where
$\colon\!\phi^0\!\colon\,=1$ and $\colon\!\phi^1\!\colon\,=\phi$.
\end{lemma}

\begin{proof}
At the regularised level (frequency cutoff $N$), $\phi_N$ is a smooth function and $\eta$ is
smooth, so $\phi_N+\eta$ is a smooth function with pointwise variance $c_N$ (the variance
of $\phi_N$; $\eta$ is deterministic and contributes no variance, so the natural Wick
constant for $(\phi_N+\eta)^n$ relative to $\mu_0$ is again $c_N$). The Hermite polynomials
satisfy the shift identity
\begin{equation}\label{eq:Hshift}
H_n(x+y;c) \;=\; \sum_{k=0}^n\binom{n}{k}y^{n-k}H_k(x;c),\qquad x,y\in\mathbb{R},\ c\ge 0.
\end{equation}
Indeed, the generating identity $\sum_n t^n H_n(z;c)/n!=\exp(tz-t^2c/2)$ gives
$\exp(t(x+y)-t^2c/2)=\exp(ty)\cdot\exp(tx-t^2c/2)$, i.e.
$\sum_n t^n H_n(x+y;c)/n!=\bigl(\sum_m t^m y^m/m!\bigr)\bigl(\sum_k t^k H_k(x;c)/k!\bigr)$,
and matching the coefficient of $t^n$ yields \eqref{eq:Hshift}. Setting
$x=\phi_N(z)$, $y=\eta(z)$, $c=c_N$ and noting $\colon\!\phi_N^k(z)\!\colon=H_k(\phi_N(z);c_N)$
from \eqref{eq:Wick}, we get pointwise:
\[
\colon\!(\phi_N+\eta)^n(z)\!\colon \;=\; \sum_{k=0}^n\binom{n}{k}\eta(z)^{n-k}\colon\!\phi_N^k(z)\!\colon.
\]
Now take $N\to\infty$. By the standard construction of Wick powers \cite[Ch.~I, \S3]{Simon},
$\colon\!\phi_N^k\!\colon\to\colon\!\phi^k\!\colon$ as $N\to\infty$ in
$L^p(\mu_0;\mathcal{C}^{-(k-1)/2-\kappa})$ for every $p<\infty$ and $\kappa>0$. The smooth
deterministic factors $\eta^{n-k}\in L^\infty(\mathbb{T}^3)$ (compact torus, smooth $\eta$),
so multiplication by $\eta^{n-k}$ is continuous on the Besov spaces involved, and we may
pass to the limit term-by-term in $L^p$.
\end{proof}

\begin{lemma}[Shift of the renormalised potential]\label{lem:Vshift}
For every $\eta\in C^\infty(\mathbb{T}^3)$,
\begin{equation}\label{eq:Vshift}
V(\phi+\eta)-V(\phi) \;=\; \int_{\mathbb{T}^3}\bigl(F_3(\eta)\colon\!\phi^3\!\colon
+ F_2(\eta)\colon\!\phi^2\!\colon + F_1(\eta)\,\phi\bigr)\,dx \;+\; F_0(\eta),
\end{equation}
in $L^p(\mu_0)$ for every $p\in[1,\infty)$, where
\begin{align}
F_3(\eta)(x) &= \lambda\,\eta(x), \label{eq:F3}\\
F_2(\eta)(x) &= \tfrac{3\lambda}{2}\eta(x)^2 + a, \label{eq:F2}\\
F_1(\eta)(x) &= \lambda\,\eta(x)^3 + 2a\,\eta(x), \label{eq:F1}\\
F_0(\eta) &= \int_{\mathbb{T}^3}\!\Bigl(\tfrac{\lambda}{4}\eta(x)^4 + a\,\eta(x)^2\Bigr)\,dx. \label{eq:F0}
\end{align}
The additive constant $b$ in \eqref{eq:V} cancels in the difference.
\end{lemma}

\begin{proof}
Apply Lemma~\ref{lem:Wickshift} (with the choice $\eta\to\eta$) to $n=4$ and $n=2$:
\begin{align*}
\colon\!(\phi+\eta)^4\!\colon &= \colon\!\phi^4\!\colon + 4\eta\colon\!\phi^3\!\colon
+ 6\eta^2\colon\!\phi^2\!\colon + 4\eta^3\phi + \eta^4,\\
\colon\!(\phi+\eta)^2\!\colon &= \colon\!\phi^2\!\colon + 2\eta\,\phi + \eta^2.
\end{align*}
Substitute into $V(\phi+\eta)$ from \eqref{eq:V}:
\begin{align*}
V(\phi+\eta)-V(\phi)
&= \tfrac{\lambda}{4}\!\int\!\bigl(4\eta\colon\!\phi^3\!\colon + 6\eta^2\colon\!\phi^2\!\colon + 4\eta^3\phi + \eta^4\bigr)dx
\;+\; a\!\int\!\bigl(2\eta\phi + \eta^2\bigr)dx \\
&= \int\!\bigl(\lambda\eta\colon\!\phi^3\!\colon
+ (\tfrac{3\lambda}{2}\eta^2+a)\colon\!\phi^2\!\colon
+ (\lambda\eta^3+2a\eta)\phi\bigr)dx
\;+\; \int\!\bigl(\tfrac{\lambda}{4}\eta^4+a\eta^2\bigr)dx,
\end{align*}
which is exactly \eqref{eq:Vshift}--\eqref{eq:F0}. The $b$ term cancels because it is the
same deterministic constant in $V(\phi+\eta)$ and $V(\phi)$. The $L^p(\mu_0)$ identity is
inherited from \eqref{eq:V} and Lemma~\ref{lem:Wickshift}, both of which converge in every
$L^p(\mu_0)$.
\end{proof}

\section{$L^p$ Bounds for the Wick-Polynomial Difference}

For any $\eta\in C^\infty(\mathbb{T}^3)$ define
\begin{equation}\label{eq:Wpsi}
W_\eta(\phi) \;:=\; V(\phi+\eta)-V(\phi)
\;=\; W_\eta^{(3)}(\phi) + W_\eta^{(2)}(\phi) + W_\eta^{(1)}(\phi) + W_\eta^{(0)},
\end{equation}
\begin{equation}\label{eq:Wj}
W_\eta^{(j)}(\phi) := \int_{\mathbb{T}^3} F_j(\eta)\colon\!\phi^j\!\colon\,dx\quad(j=1,2,3),
\qquad W_\eta^{(0)} := F_0(\eta).
\end{equation}
Here $W_\eta^{(0)}$ is a finite deterministic constant (since $F_0(\eta)$ is the integral of
a smooth function over the compact torus).

\begin{lemma}\label{lem:WickL2}
For each $j\in\{1,2,3\}$, $W_\eta^{(j)}\in L^2(\mu_0)$, with explicit second moments
\begin{align}
\mathbb{E}[(W_\eta^{(1)})^2] &= \langle F_1(\eta),\,C\,F_1(\eta)\rangle_{L^2},\label{eq:m1}\\
\mathbb{E}[(W_\eta^{(2)})^2] &= 2\!\int_{\mathbb{T}^3\times\mathbb{T}^3}\! F_2(\eta)(x)F_2(\eta)(y)\,C(x,y)^2\,dx\,dy,\label{eq:m2}\\
\mathbb{E}[(W_\eta^{(3)})^2] &= 6\!\int_{\mathbb{T}^3\times\mathbb{T}^3}\! F_3(\eta)(x)F_3(\eta)(y)\,C(x,y)^3\,dx\,dy,\label{eq:m3}
\end{align}
all of which are finite. Here $C(x,y)$ is the integral kernel of $C=(-\Delta+m^2)^{-1}$ on
$\mathbb{T}^3$.
\end{lemma}

\begin{proof}
The formulae \eqref{eq:m1}, \eqref{eq:m2}, \eqref{eq:m3} follow from the standard identity
$\mathbb{E}[\colon\!\phi^j(x)\!\colon\colon\!\phi^k(y)\!\colon]=\delta_{jk}\,j!\,C(x,y)^j$
(the $L^2$-orthogonality of distinct Wiener chaos and the explicit $j$-th-chaos covariance,
\cite[Theorem~I.3]{Simon}), pairing against the smooth test functions $F_j(\eta)$.
We must verify finiteness for $j=1,2,3$.

\emph{$j=1$:} $F_1(\eta)\in C^\infty(\mathbb{T}^3)\subset L^2(\mathbb{T}^3)$, $C$ is bounded
on $L^2$, so \eqref{eq:m1} is finite.

\emph{$j=2$:} The torus heat-kernel/Green-function bound
$|C(x,y)|\le c_0|x-y|^{-1}+c_1$ for $x\ne y$ on $\mathbb{T}^3$ (standard short-distance
expansion of $(-\Delta+m^2)^{-1}$ in three dimensions, e.g.\ \cite[Lemma~3.4]{GlimmJaffe})
gives $C(x,y)^2\le 2c_0^2|x-y|^{-2}+2c_1^2$ for $x\ne y$. Since $|x-y|^{-2}$ is locally
integrable in $\mathbb{R}^3$ (the volume element is $r^2\,dr$ in spherical coordinates,
making $\int_0^1 r^{-2}\cdot r^2\,dr=1$ finite, and the bound transfers to the torus by
periodicity), the double integral in \eqref{eq:m2} is finite for the bounded function
$F_2(\eta)$.

\emph{$j=3$:} We compute \eqref{eq:m3} in Fourier. The kernel $C(x,y)$ is translation
invariant: $C(x,y)=K(x-y)$ with $\widehat K(k)=(4\pi^2|k|^2+m^2)^{-1}$. By Parseval and
convolution,
\begin{equation}\label{eq:m3-Fourier}
\mathbb{E}[(W_\eta^{(3)})^2]
\;=\; 6\sum_{k\in\mathbb{Z}^3}|\widehat{F_3(\eta)}(k)|^2\,M_3(k),
\quad
M_3(k):=\sum_{\substack{k_1,k_2\in\mathbb{Z}^3 \\ k_1+k_2+k_3=k}}\,
\prod_{i=1}^{3}\!\frac{1}{4\pi^2|k_i|^2+m^2},
\end{equation}
where the sum is over $k_1,k_2\in\mathbb{Z}^3$ and $k_3:=k-k_1-k_2$. We bound $M_3(k)$.
Use the elementary inequality $(4\pi^2|k_i|^2+m^2)^{-1}\le c_*\langle k_i\rangle^{-2}$ with
$\langle k\rangle:=(1+|k|^2)^{1/2}$ and $c_*:=\max(1/m^2, 1/(4\pi^2))$. Hence
\[
M_3(k)\;\le\; c_*^3\!\sum_{k_1,k_2}\langle k_1\rangle^{-2}\langle k_2\rangle^{-2}\langle k-k_1-k_2\rangle^{-2}.
\]
For fixed $k$, the inner sum is uniformly bounded by a constant $C_3$ independent of $k$:
indeed, the discrete convolution $\langle\cdot\rangle^{-2}*\langle\cdot\rangle^{-2}$ on
$\mathbb{Z}^3$ is bounded pointwise by $c'\langle\cdot\rangle^{-1}$ (this is the analogue
in $\mathbb{Z}^3$ of $\int_{\mathbb{R}^3}\langle\xi\rangle^{-2}\langle\xi-z\rangle^{-2}d\xi
\le c\langle z\rangle^{-1}$, valid because
$2+2-3=1$ controls the decay), and a second convolution with $\langle\cdot\rangle^{-2}$
gives a bounded function on $\mathbb{Z}^3$ (since
$\sum_z \langle z\rangle^{-1}\langle k-z\rangle^{-2}\le c''$, finite uniformly in $k$ by the
same Hardy--Littlewood--Sobolev/discrete-convolution argument with exponents $1+2-3=0$, the
critical case for boundedness). Thus $M_3(k)\le c_*^3\cdot c'\cdot c''=:C_3<\infty$
uniformly in $k\in\mathbb{Z}^3$. Substituting,
\[
\mathbb{E}[(W_\eta^{(3)})^2]\;\le\;6C_3\sum_{k\in\mathbb{Z}^3}|\widehat{F_3(\eta)}(k)|^2
\;=\;6C_3\,\|F_3(\eta)\|_{L^2(\mathbb{T}^3)}^2\;<\;\infty,
\]
using Parseval and $F_3(\eta)=\lambda\eta\in L^2$.
\end{proof}

\begin{lemma}\label{lem:WickLp}
For each $j\in\{1,2,3\}$ and every $p\in[1,\infty)$, $W_\eta^{(j)}\in L^p(\mu_0)$.
Consequently $W_\eta\in L^p(\mu_0)$ for every $p\in[1,\infty)$.
\end{lemma}

\begin{proof}
Each $W_\eta^{(j)}$ ($j=1,2,3$) belongs to the $j$-th Wiener chaos $\mathcal{H}_j$ of
$\mu_0$: this is by construction, since $W_\eta^{(j)}$ is the $L^2$-limit (as $N\to\infty$)
of $\int F_j(\eta)\colon\!\phi_N^j\!\colon dx$, and each cutoff is a $j$-th-degree
polynomial in jointly Gaussian variables which equals the projection onto $\mathcal{H}_j$
by Wick orthogonality \cite[Ch.~I]{Simon}. By Borell's hypercontractivity bound on Wiener
chaos \cite{Borell}: for any $X\in\mathcal{H}_j$ and any $p\ge 2$,
\begin{equation}\label{eq:Borell}
\|X\|_{L^p(\mu_0)}\;\le\;(p-1)^{j/2}\,\|X\|_{L^2(\mu_0)}.
\end{equation}
Applied with $j=1,2,3$ and the $L^2$-bound from Lemma~\ref{lem:WickL2}, this gives
$\|W_\eta^{(j)}\|_{L^p}<\infty$ for every $p\ge 2$, hence (by Jensen) for every $p\ge 1$.
Adding the deterministic constant $W_\eta^{(0)}$ and using Minkowski,
$W_\eta\in L^p(\mu_0)$.
\end{proof}

\section{Exponential Integrability of the Combined Density}

We now show that the random variable that will play the role of the Radon--Nikodym
derivative is positive and integrable. Write
\begin{equation}\label{eq:Geta}
G_\eta(\phi) \;:=\; e^{-W_\eta(\phi)}\,R_\eta(\phi),
\qquad \eta\in C^\infty(\mathbb{T}^3).
\end{equation}

\begin{lemma}[$L^p(\mu_0)$-integrability of $G_\eta$]\label{lem:GLp}
For every $\eta\in C^\infty(\mathbb{T}^3)$ and every $p\in[1,\infty)$,
$G_\eta\in L^p(\mu_0)$ and $G_\eta>0$ $\mu_0$-a.s.
\end{lemma}

\begin{proof}
Positivity of $G_\eta$ follows from $R_\eta>0$ everywhere (it is an exponential of a real
Gaussian, finite $\mu_0$-a.s.) and $W_\eta\in\mathbb{R}$ $\mu_0$-a.s.\ (Lemma~\ref{lem:WickLp}).

For integrability we use the following identity, which is the key analytic step. By
Lemma~\ref{lem:CMfree} and the change-of-variables formula,
\begin{equation}\label{eq:key-identity}
\mathbb{E}_{\mu_0}\!\bigl[e^{p(-W_\eta)}\,R_\eta^p\bigr]
\;=\; \mathbb{E}_{\mu_0}\!\bigl[e^{-pW_\eta(\phi)}R_\eta(\phi)^{p}\bigr].
\end{equation}
We bound the right-hand side directly. By the definition $W_\eta=V(\phi+\eta)-V(\phi)$
(equation \eqref{eq:Wpsi}),
\[
e^{-pW_\eta(\phi)}\,R_\eta(\phi)^{p}
\;=\; e^{-p\,V(\phi+\eta)}\,e^{p\,V(\phi)}\,R_\eta(\phi)^{p}.
\]
H\"older's inequality with three factors $\frac{1}{r}+\frac{1}{s}+\frac{1}{t}=1$,
$r,s,t>1$, gives
\begin{equation}\label{eq:holder}
\mathbb{E}_{\mu_0}[G_\eta^p]
\;\le\; \bigl\|e^{-pV(\phi+\eta)}\bigr\|_{L^r(\mu_0)}\cdot
\bigl\|e^{pV(\phi)}\bigr\|_{L^s(\mu_0)}\cdot
\bigl\|R_\eta^p\bigr\|_{L^t(\mu_0)}.
\end{equation}
Each of the three factors is finite:
\begin{itemize}
\item[(i)] The second factor: $e^{pV}\in L^s(\mu_0)$ for every $s<\infty$ by
Theorem~\ref{thm:GJ} (two-sided $L^p$-bound on $e^{\pm V}$).
\item[(ii)] The third factor: $R_\eta^p\in L^t(\mu_0)$ for every $t<\infty$ by
Lemma~\ref{lem:Lp-Rpsi}.
\item[(iii)] The first factor: $\bigl\|e^{-pV(\phi+\eta)}\bigr\|_{L^r(\mu_0)}^r
=\mathbb{E}_{\mu_0}[e^{-prV(\phi+\eta)}]
=\int e^{-prV(\phi)}\,d(T_\eta^*\mu_0)(\phi)
=\int e^{-prV(\phi)}\,R_\eta(\phi)\,d\mu_0(\phi)$
(using Lemma~\ref{lem:CMfree} again). By H\"older with $\frac{1}{u}+\frac{1}{v}=1$,
$u,v>1$, the right-hand side is bounded by
$\|e^{-prV}\|_{L^u(\mu_0)}\|R_\eta\|_{L^v(\mu_0)}$, which is finite by
Theorem~\ref{thm:GJ} and Lemma~\ref{lem:Lp-Rpsi}.
\end{itemize}
Hence $\mathbb{E}_{\mu_0}[G_\eta^p]<\infty$, i.e.\ $G_\eta\in L^p(\mu_0)$.
\end{proof}

\begin{lemma}\label{lem:Ginv-Lp}
For every $\eta\in C^\infty(\mathbb{T}^3)$ and every $p\in[1,\infty)$,
$1/G_\eta=e^{W_\eta}/R_\eta\in L^p(\mu_0)$.
\end{lemma}

\begin{proof}
By the same H\"older decomposition as in \eqref{eq:holder}, with sign of $W_\eta$ and
$\log R_\eta$ flipped:
\[
\bigl(e^{W_\eta}/R_\eta\bigr)^p \;=\; e^{p\,V(\phi+\eta)}\,e^{-p\,V(\phi)}\,R_\eta^{-p}.
\]
By H\"older,
\[
\mathbb{E}_{\mu_0}[(1/G_\eta)^p]
\le\bigl\|e^{pV(\phi+\eta)}\bigr\|_{L^r}\bigl\|e^{-pV}\bigr\|_{L^s}\bigl\|R_\eta^{-p}\bigr\|_{L^t},
\]
each factor finite by Theorem~\ref{thm:GJ}, Theorem~\ref{thm:GJ} (two-sided bound), and
Lemma~\ref{lem:Lp-Rpsi} (with $\alpha=-p$) respectively, the first factor handled by the
Cameron--Martin substitution as in step (iii) of Lemma~\ref{lem:GLp}.
\end{proof}

\section{Equivalence of $\mu$ and $T_\psi^*\mu$}

We now combine the pieces. We work directly with the convention
$\eta=-\psi$, which makes the resulting Radon--Nikodym derivative most transparent.

\begin{theorem}[Main result]\label{thm:main}
Let $\psi\in C^\infty(\mathbb{T}^3)$ be smooth (not necessarily nonzero). Then $\mu$ and
$T_\psi^*\mu$ are mutually absolutely continuous, and the Radon--Nikodym derivative is
\begin{equation}\label{eq:RN-mu}
\frac{d(T_\psi^*\mu)}{d\mu}(\phi) \;=\; e^{-W_{-\psi}(\phi)}\,R_\psi(\phi)
\;=\; \exp\!\bigl(V(\phi)-V(\phi-\psi)\bigr)\,R_\psi(\phi),
\end{equation}
where $W_{-\psi}(\phi)=V(\phi-\psi)-V(\phi)$ as in \eqref{eq:Wpsi}.
\end{theorem}

\begin{proof}
\textbf{Step 1: pushforward identity.}
For any bounded measurable $g:\mathcal{D}'(\mathbb{T}^3)\to\mathbb{R}$,
\[
\int g\,d(T_\psi^*\mu)\;=\;\int g(\phi+\psi)\,d\mu(\phi)
\;=\;\frac{1}{Z}\int g(\phi+\psi)\,e^{-V(\phi)}\,d\mu_0(\phi),
\]
the second equality by Theorem~\ref{thm:GJ}.

\textbf{Step 2: Cameron--Martin substitution.}
The map $\phi\mapsto\phi+\psi$ pushes $\mu_0$ to $T_\psi^*\mu_0$, which by
Lemma~\ref{lem:CMfree} has density $R_\psi$ w.r.t.\ $\mu_0$. Substituting $\phi':=\phi+\psi$
(so $\phi=\phi'-\psi$):
\begin{align*}
\frac{1}{Z}\int g(\phi+\psi)\,e^{-V(\phi)}\,d\mu_0(\phi)
&= \frac{1}{Z}\int g(\phi')\,e^{-V(\phi'-\psi)}\,d(T_\psi^*\mu_0)(\phi')\\
&= \frac{1}{Z}\int g(\phi')\,e^{-V(\phi'-\psi)}\,R_\psi(\phi')\,d\mu_0(\phi').
\end{align*}

\textbf{Step 3: rewrite using $W_{-\psi}$.}
By Lemma~\ref{lem:Vshift} applied with $\eta=-\psi$,
\[
V(\phi'-\psi)-V(\phi') \;=\; W_{-\psi}(\phi')\quad\text{in } L^p(\mu_0)\ \forall p<\infty,
\]
so $e^{-V(\phi'-\psi)}=e^{-V(\phi')}e^{-W_{-\psi}(\phi')}$. Substituting,
\begin{align*}
\int g\,d(T_\psi^*\mu)
&= \frac{1}{Z}\int g(\phi')\,e^{-V(\phi')}\,e^{-W_{-\psi}(\phi')}\,R_\psi(\phi')\,d\mu_0(\phi')\\
&= \int g(\phi')\,e^{-W_{-\psi}(\phi')}\,R_\psi(\phi')\,d\mu(\phi'),
\end{align*}
using $d\mu=Z^{-1}e^{-V}d\mu_0$.

\textbf{Step 4: identify the density.}
Take $g\equiv 1$: the left-hand side is $1$, so
$\int e^{-W_{-\psi}}\,R_\psi\,d\mu=1$, in particular $e^{-W_{-\psi}}R_\psi\in L^1(\mu)$.
For general $g$ the identity from Step~3 reads
\[
\int g\,d(T_\psi^*\mu)\;=\;\int g\cdot \bigl(e^{-W_{-\psi}}R_\psi\bigr)\,d\mu,
\]
which is exactly the statement that $e^{-W_{-\psi}}R_\psi$ is the Radon--Nikodym derivative
$d(T_\psi^*\mu)/d\mu$. This proves $T_\psi^*\mu\ll\mu$ and gives \eqref{eq:RN-mu}.

\textbf{Step 5: $\mu$-integrability of the density.}
The function $e^{-W_{-\psi}}R_\psi$ is positive $\mu$-a.s.\ since $\mu\ll\mu_0$ and the
function is positive $\mu_0$-a.s. (Lemma~\ref{lem:GLp} applied with $\eta=-\psi$, noting
$R_\psi=R_{-(-\psi)}$ does not affect positivity). It is in $L^1(\mu)$ by Step~4 (the
integral equals $1$).

\textbf{Step 6: $\mu\ll T_\psi^*\mu$.}
A measure $\nu_1$ is absolutely continuous w.r.t.\ $\nu_2$ with strictly positive density
$\rho$ iff $\nu_2$ is absolutely continuous w.r.t.\ $\nu_1$ with density $1/\rho$. The
density $\rho:=e^{-W_{-\psi}}R_\psi$ is strictly positive $\mu$-a.s.\ (Step 5), so
$\mu\ll T_\psi^*\mu$ with density
\[
\frac{d\mu}{d(T_\psi^*\mu)}(\phi) \;=\; e^{W_{-\psi}(\phi)}/R_\psi(\phi).
\]
This density is in $L^1(T_\psi^*\mu)$: indeed,
\[
\int e^{W_{-\psi}}/R_\psi\,d(T_\psi^*\mu)
\;=\; \int e^{W_{-\psi}}/R_\psi\cdot e^{-W_{-\psi}}R_\psi\,d\mu
\;=\; \int 1\,d\mu \;=\; 1,
\]
which confirms (equivalently) the consistency of the two densities.

Combining Steps 1--6 yields $\mu\sim T_\psi^*\mu$ with the Radon--Nikodym derivative
\eqref{eq:RN-mu}.
\end{proof}

\section{Boundary and Degenerate Cases}

The problem specifies $\psi\not\equiv 0$ and smooth, but the argument works for every
smooth $\psi$:

\begin{itemize}
\item If $\psi\equiv0$, then $T_\psi=\mathrm{id}$ and $T_\psi^*\mu=\mu$ trivially. Indeed
$R_0=\exp(0-0)=1$ and $W_0=V(\phi)-V(\phi)=0$, so the density in \eqref{eq:RN-mu} is $1$,
and the claim is the trivial identity.
\item For every smooth $\psi\ne 0$, Theorem~\ref{thm:main} gives equivalence with the
non-trivial density \eqref{eq:RN-mu}.
\item For $\psi\notin\mathcal{H}=H^1(\mathbb{T}^3)$, the Cameron--Martin theorem
(Theorem~\ref{thm:CM}) yields $T_\psi^*\mu_0\perp\mu_0$. Since $\mu\sim\mu_0$ (the density
$e^{-V}/Z$ is strictly positive $\mu_0$-a.s.\ and $\mu_0$-integrable), $T_\psi^*\mu$ is
also singular w.r.t.\ $\mu$, so the
analogous question would have the opposite answer (\emph{No}). This case is excluded by the
problem's hypothesis $\psi\in C^\infty\subset H^1$.
\end{itemize}

\section{Final Statement}

\begin{theorem}[Restatement]\label{thm:final}
Let $\mu$ be the $\Phi^4_3$ measure on $\mathcal{D}'(\mathbb{T}^3)$, with the
representation $d\mu=Z^{-1}e^{-V}d\mu_0$ from Theorem~\ref{thm:GJ}, and let
$\psi\in C^\infty(\mathbb{T}^3)$ be smooth and not identically zero. Then $\mu$ and
$T_\psi^*\mu$ are equivalent, with explicit Radon--Nikodym derivative
\begin{equation}\label{eq:final-RN}
\frac{d(T_\psi^*\mu)}{d\mu}(\phi)
\;=\; \exp\!\Bigl(\,V(\phi)-V(\phi-\psi)
\;+\; \langle\phi,(-\Delta+m^2)\psi\rangle
\;-\; \tfrac12\langle\psi,(-\Delta+m^2)\psi\rangle\Bigr).
\end{equation}
The renormalised potential difference $V(\phi)-V(\phi-\psi)=-W_{-\psi}(\phi)$ is the
$L^p(\mu_0)$-defined random variable given (with $\eta=-\psi$ in
Lemma~\ref{lem:Vshift}) by
\[
V(\phi)-V(\phi-\psi)
\;=\; -\int\!\bigl(F_3(-\psi)\colon\!\phi^3\!\colon
+ F_2(-\psi)\colon\!\phi^2\!\colon
+ F_1(-\psi)\,\phi\bigr)dx
\;-\; F_0(-\psi),
\]
i.e.\ explicitly
\[
V(\phi)-V(\phi-\psi)
\;=\; \int\!\Bigl(\lambda\psi\colon\!\phi^3\!\colon
\;-\; \bigl(\tfrac{3\lambda}{2}\psi^2+a\bigr)\colon\!\phi^2\!\colon
\;+\; \bigl(\lambda\psi^3+2a\psi\bigr)\phi\Bigr)dx
\;-\; \int\!\bigl(\tfrac{\lambda}{4}\psi^4 + a\psi^2\bigr)dx,
\]
using the parities $F_3(-\psi)=-\lambda\psi$, $F_2(-\psi)=\tfrac{3\lambda}{2}\psi^2+a$,
$F_1(-\psi)=-(\lambda\psi^3+2a\psi)$, $F_0(-\psi)=\int(\tfrac{\lambda}{4}\psi^4+a\psi^2)dx$.
The linear form $\langle\phi,(-\Delta+m^2)\psi\rangle$ is the centred Gaussian variable on
$(\mathcal{D}'(\mathbb{T}^3),\mu_0)$ with variance
$\langle\psi,(-\Delta+m^2)\psi\rangle$.
\end{theorem}

\begin{proof}
This is Theorem~\ref{thm:main} together with \eqref{eq:RPsi}, the explicit substitution
$\eta\mapsto-\psi$ in Lemma~\ref{lem:Vshift}, and the cancellation $-(-W_{-\psi})=W_{-\psi}$
applied with the sign that yields the final form \eqref{eq:final-RN}. The strict positivity
of the right-hand side is immediate (it is an exponential of a real number).
\end{proof}

\paragraph{Conclusion.}
The answer to the problem is \textbf{Yes}: $\mu$ and $T_\psi^*\mu$ are equivalent for
every smooth $\psi$ on $\mathbb{T}^3$ (including the case $\psi\not\equiv 0$ specified in
the problem statement), with the explicit, strictly positive Radon--Nikodym derivative
\eqref{eq:final-RN}.

\paragraph{References.}
\begin{itemize}
\item[\cite{Bogachev}] V.~I.~Bogachev, \emph{Gaussian Measures}, Math.\ Surveys
Monographs~62, AMS, 1998 (Theorem~2.4.5: Cameron--Martin).
\item[\cite{Borell}] C.~Borell, ``Tail probabilities in Gauss space,'' in
\emph{Vector space measures and applications I}, LNM~644, Springer, 1978 (hypercontractive
bounds on Wiener chaos).
\item[\cite{CatellierChouk}] R.~Catellier, K.~Chouk, ``Paracontrolled distributions and
the 3-dimensional stochastic quantization equation,'' \emph{Annals of Probability} 46(5),
2018.
\item[\cite{GH}] M.~Gubinelli, M.~Hofmanov\'a, ``A PDE construction of the Euclidean
$\Phi^4_3$ quantum field theory,'' \emph{Comm.\ Math.\ Phys.} 384, 2021.
\item[\cite{GlimmJaffe}] J.~Glimm, A.~Jaffe, \emph{Quantum Physics: A Functional Integral
Point of View}, 2nd edition, Springer, 1987 (existence of $\Phi^4_3$, two-sided
$L^p$-bounds on $e^{\pm V}$).
\item[\cite{HairerRS}] M.~Hairer, ``A theory of regularity structures,''
\emph{Inventiones Math.} 198, 2014.
\item[\cite{Simon}] B.~Simon, \emph{The $P(\phi)_2$ Euclidean (quantum) field theory},
Princeton University Press, 1974 (Wick power $L^p$-convergence and orthogonality of
chaos).
\end{itemize}

\end{document}
